\section{Developing the HistoryFinder Tool and Algorithm}
\label{sec:history-finder-algorithm}

The first, third, and fourth authors collaboratively designed the \textit{HistoryFinder} algorithm through multiple Zoom meetings. The first author has 8+ years of industry experience; the third has 14+ years in industry and 3+ in research; and the last has 1 year in industry and 15 years in research. During these sessions, they analyzed the strengths and limitations of two existing approaches, \textit{CodeShovel} and \textit{CodeTracker}, which guided the design of their new algorithm. Their analysis revealed that while \textit{CodeTracker} is generally more accurate than \textit{CodeShovel}, it suffers from significantly slower runtime performance. This is primarily due to its dependence on RefactorMiner, whereas \textit{CodeShovel} uses a simpler, faster string-matching algorithm. To balance accuracy and efficiency, \textit{HistoryFinder} adopts \textit{CodeShovel’s} lightweight string-matching approach but improves upon it with a more precise method-tracking strategy. The main algorithmic differences are as follows.
\begin{itemize}
    \item In Phase 2 of the \textit{CodeShovel} algorithm, the first method with body similarity $\ge0.75$ is selected, without checking for a matching signature. This can lead to tracking an extracted method instead of the original. In contrast, \textit{HistoryFinder} first checks for a matching signature and resorts to body similarity only if none is found.
\item When the method is not found in the same file, \textit{CodeShovel} performs two phases: first, it looks for methods with the same signature (even with low body similarity), then for methods with high body similarity. During the training phase, a method in a different file with the same signature but low similarity was observed to not necessarily be the same method. Merging these phases---directly searching for highly similar methods without checking the signature---improves precision.

\item Additionally, in Phase 4, \textit{CodeShovel} compares all methods across modified files, which is not only time-consuming but can also lead to incorrect matches, especially in the case of method overloading. This issue was addressed by filtering out methods that were unchanged in both commits, as such methods cannot be the ones of interest. Also, \textit{CodeShovel} returns the first score exceeding the threshold, which is problematic when a more similar method exists. To alleviate this issue, \textit{HistoryFinder} returns the best matching method passing the similarity threshold.
\end{itemize}

\textit{HistoryFinder} also improved some of the technical weaknesses that \textit{CodeShovel} has. For example, during development, if code is merged without rebasing, conflicts can arise when a method is modified concurrently in different branches. The resolution of these conflicts may introduce changes to the method’s history that must be considered. \textit{CodeShovel}, however, excludes merge commits and ignores changes introduced during merges. In this approach, merge commits are included, and the method’s state is compared with each parent commit, allowing accurate detection of any changes introduced from either parent branch. Moreover, \textit{CodeShovel} does not account for formatting changes, annotation, or JavaDoc updates, whereas \textit{HistoryFinder} does. Another limitation of \textit{CodeShovel} is its case-insensitive string-matching algorithm. This approach can be problematic, especially when comparing quoted string literals within a method body, as it may overlook subtle but important differences that could influence the method's behavior.

\subsection{Overview of HistoryFinder}
The \textit{HistoryFinder} tool needs several key inputs to trace the history of a method. These include the \texttt{Git} repository URL or local directory, a specific commit hash (SHA), the file path of the method (including the file name), the name of the method, and the line number where the method declaration begins. The output is a sequential list of commit hashes, each accompanied by relevant metadata. This metadata provides granularity of changes, such as changes to the method signature, body, and annotations, and identifies the contributor responsible for each modification. When executed via the command line, the output is saved as a \texttt{JSON} file.

The implementation utilizes the \texttt{JGit} library\footnote{https://github.com/eclipse-jgit/jgit last accessed: Jul 01, 2025}—a \texttt{Java}-based \texttt{Git} implementation—for repository cloning, commit traversal, and fine-grained file operations, including detection of added, deleted, moved, and modified files across revisions, and employs \texttt{JavaParser}\footnote{https://github.com/javaparser/javaparser last accessed: Jul 01, 2025} to generate an Abstract Syntax Tree (AST). Once the AST is constructed, it allows traversal through the list of methods in the file and extraction of specific components, such as the method signature, body, annotation, and JavaDoc. For method comparison, the implementation uses the Jaro-Winkler edit distance~\cite{winkler90} score similarity, as in \textit{CodeShovel}, between method signatures and bodies. Although \texttt{SimHash}~\cite{charikar_similarity_2002} was also tested for similarity matching, Jaro-Winkler consistently provides better recall than \texttt{SimHash}.

When the commit history is interpreted as a hierarchical graph structure, all commits leading up to a specific point in a \texttt{Git} project form a Directed Acyclic Graph (DAG)~\cite{geisshirt_git_2018}, representing the project's complete commit history to that point. To identify all commits that modified a method up to a given version (e.g., the starting commit), the DAG is traversed backward—from the target commit to its ancestors—until the commit where the method was originally introduced is reached. During this traversal, if a commit has a single parent, the method in the current commit is compared with its counterpart in the parent commit to detect changes. In the case of merge commits with multiple parents, the method may exist in any of the parent commits. Given that a project may contain thousands of commits (as shown in Table~\ref{tab:history-finder-repository-list}), traversing the entire DAG can be time-consuming. To optimize this process, \textit{HistoryFinder} leverages the \texttt{git log <commit> <file>} command, which filters the commit history to only those that modified the file containing the target method. This significantly reduces the search space, allowing traversal of a lightweight, file-specific DAG to detect method changes efficiently. However, if the method has been moved to a different file or the file itself has been renamed, a new DAG must be reconstructed from that point onward using the new file path.

\subsection{The HistoryFinder Algorithm}
Algorithm~\ref{alg:algorithm-history-finder} outlines how the file-specific DAG is constructed and traversed to detect the change history of a method. The procedure \textit{HistoryFinder} is initially invoked with the starting commit, file path, method name, and its starting line number. It may be invoked recursively if the method is detected to have moved to a different file or if the file itself has been renamed. To build the DAG, the utility procedure \textit{GITLOG} is first called (line 2), which internally executes the \texttt{git log <commit> <file>} command. This retrieves only the commits that have modified the target file containing the method. Using these filtered commits, a directed acyclic graph \( G \) is constructed, where each node represents a commit that changed the file. In the DAG, an edge is drawn from node $u$ to node $v$ if the file is modified in commit $u$, followed by in commit $v$, establishing a direct connection between them. Once the DAG is constructed, it is traversed using the Breadth-First Search (BFS) algorithm~\cite{cormen_introduction_2022} to systematically explore all relevant commits.

Queue \( Q \) keeps track of all discovered commit nodes that need to be explored. To unambiguously locate the target method in a commit, each commit node is represented by a structure containing the commit hash, file path, method name, and line number, created by invoking \texttt{Create} (line 3 for the initial node) procedure. The queue \( Q \) is initialized with the starting node (line 4), which is also the only entry in the set \texttt{visited} to mark it as explored. Whenever a change to the method is detected, it is recorded in the history list \( H \), which is initially empty (line 6). The \texttt{while} loop iterates through all discovered nodes. In each iteration, an unexplored commit node \( u \) is dequeued from the head of \( Q \) (line 8). The method already found in commit node \( u \) is located by calling the procedure \texttt{Find} (line 9), which takes the commit hash, file, method name, and line number as parameters, accessed using dot notation. Next, the parent node \( p \) of \( u \) is retrieved (line 10), where the method might be found. From lines 11 to 18, the algorithm checks whether the method exists in the direct parent commit \( p \), accounting for possible modifications.

\begin{algorithm}
\caption{HistoryFinder Algorithm}
\label{alg:algorithm-history-finder}
\footnotesize
\SetKwFunction{hf}{HistoryFinder}
\SetKwProg{Procedure}{procedure}{:}{end procedure}

\Procedure{\hf{$hash,file,name,line$}}{ %\tcp*[r]{$name$ and $line$ are the method's name and start line}
    
    \KwIn{
      $hash$: Start commit hash\linebreak
      \Indp
        $file$: File containing the method\linebreak
        $name$: Method name\linebreak
        $line$: Start line of the method
      \Indm
    }
    \KwOut{Change history of the input method}
    $G \gets \textsc{GitLog}(hash, file)$\tcp*[r]{Build file-level DAG}
    $s \gets \textsc{Create}(hash,file,name,line)$\;
    $Q \gets \{s\}$\tcp*[r]{Initialize queue with start node}
    
    $visited \gets \{s\}$\tcp*[r]{Track visited nodes}
    $H \gets \emptyset$\tcp*[r]{Initialize with empty method history}
    \While{$Q \neq \emptyset$}{ %\tcp*[r]{Process all discovered commit nodes}
        $u \gets \textsc{Dequeue}(Q)$\;
        $um \gets \textsc{Find}(u.hash, u.file, u.name, u.line)$\;
        $p \gets \textsc{Parent}(G, u)$ \tcp*[r]{Direct parent commit}
        \BlankLine
        \tcp{Search method by signature, body in original file, then body in other files}
        $m \gets \textsc{SignatureMatch}(p,file,um)$\;
        $altFile \gets \textsc{NIL}$\tcp*[r]{Track method/file move}
         \If{$m == \text{NIL}$}{ %\tcp*[r]{No match by signature}
            $m \gets \textsc{BodyMatch}(p,file,um, 0.70)$\;
            \If{$m == \text{NIL}$}{%\tcp*[r]{No match by body similarity}
                %\tcp*[r]{Examine other files}
                $m,altFile \gets \textsc{AltFileMatch}(u.hash,p,um)$;
            }
        }
        \BlankLine
        \tcp{Method found in original/other file}
        \If{$m \not= \text{NIL}$}{
            \If{$um.\text{text} \not= m.\text{text} \textbf{ or } altFile \not= \textsc{NIL}$}{
                \textsc{Add}$(H, u)$ \tcp*[r]{Record method change}
            }
            \tcp{Continue tracking in the new file}
            \If{$altFile \not= \textsc{NIL}$}{% \tcp*[r]{Method or file moved}
                $sH \gets \textsc{HistoryFinder}(m.hash,m.file,m.name,m.line)$\;
                %\vspace{-0.35cm}
                \textsc{AddAll}$(H, sH)$ \tcp*[r]{Append sub-history}
            }
        \tcp{Enqueue unvisited ancestor commits}
        \If{$m \not= \textsc{NIL} \textbf{ and } altFile == \textsc{NIL}$}{ %\tcp*[r]{Method remained in the same file}
            \ForEach{$v \in G.\text{Adj}[u]$}{
                \If{$v \notin visited$}{
                    \textsc{Add}$(visited, v)$\;
                    $v.name \gets m.name$\;
                    $v.line \gets m.line$\;
                    \textsc{Enqueue}$(Q, v)$\;
                }
            }
        }
        }
    }    
    \Return $H$\;
}
\end{algorithm}

\begin{algorithm}
\caption{Utility Procedures for Method Matching}
\label{alg:history-finder-similarity-matching}
% \small
\footnotesize
\SetKwFunction{fbsm}{SignatureMatch}
\SetKwFunction{fbbm}{BodyMatch}
\SetKwFunction{afm}{AltFileMatch}
\SetKwProg{Procedure}{procedure}{:}{end procedure}
\Procedure{\fbsm{$hash,file,tm$}}{%\tcp*[r]{Match target method by exact signature}
    \KwIn{
      $hash$: Commit hash\linebreak
      \Indp
        $file$: File containing methods\linebreak
        $tm$: Target method
      \Indm
      }
    \KwOut{Matching method or $\text{NIL}$ if not found}
    $methods \gets \textsc{ParseMethod}(hash, file)$\;%\tcp*[r]{Parse methods}
    \ForEach{$m \in methods$}{%\tcp*[r]{Iterate through parsed methods}
        %\tcp*[h]{Match signature}\;
        \If{$m.\text{signature} == tm.\text{signature}$}{
            \Return $m$ \tcp*[r]{Return matching method}
        }
    }
    \Return $\text{NIL}$ \tcp*[r]{No matching method found}
}
\BlankLine
%\tcp*[r]{Match target method by body similarity}
\Procedure{\fbbm{$hash,file, tm, threshold$}}{% \tcp*[r]{Match target method by body similarity}
    \KwIn{
      $hash$: Commit hash\linebreak
      \Indp
        $file$: File containing methods\linebreak
        $tm$: Target method\linebreak
        $threshold$: Similarity threshold for matching
      \Indm
      }
    \KwOut{
    $bm$: Best matching method meeting threshold\linebreak
      \Indp
        $bs$: Similarity score of the best match
      \Indm}
    $methods \gets \textsc{ParseMethod}(hash, file)$\; %\tcp*[r]{Parse methods}
    $bm \gets \text{NIL}$ \tcp*[r]{Initialize best matching method}
    $bs \gets 0$ \tcp*[r]{Initialize highest similarity score}
    \ForEach{$m \in methods$}{% \tcp*[r]{Iterate through all methods}
        %\tcp{Compute similarity with target method}
        $score \gets \textsc{Similarity}(m.text,tm.text)$\;
        \If{$score >= threshold \textbf{ and } score > bs$}{
            $bm \gets m$ \tcp*[r]{Update best matching method}
            $bs \gets score$ \tcp*[r]{Update highest similarity}
        }
    }
    \Return $bm,bs$\;
}
\BlankLine
\Procedure{\afm{$hash,parentHash,tm$}}{% \tcp*[r]{Search for target method in other modified files}
   \KwIn{
      $hash$: Commit hash\linebreak
      \Indp
        $parentHash$: Parent commit hash\linebreak
        $tm$: Target method
      \Indm
      }
    \KwOut{
    $bm$: Best matching method\linebreak
      \Indp
        $altFile$: File containing the best matching
      \Indm}
    %\tcp{Get files changed between commits}
    $changeFiles \gets \textsc{findChangeFiles}(hash,parentHash)$\;
    $bm \gets \text{NIL}$ \tcp*[r]{Initialize best matching method}
    $bs \gets 0$ \tcp*[r]{Initialize highest similarity score}
    $altFile \gets \textsc{NIL}$ \tcp*[r]{Best method matching file}
    \ForEach{$file \in changeFiles$}{
        $m,s \gets \textsc{BodyMatch}(parentHash,file,tm,0.75)$\;
        \If{$s > bs$}{%\tcp*[r]{Update best match if score improves}
            $bm \gets m$\;
            $bs \gets s$\;
            $altFile \gets file$\;
        }
    }
    \Return $bm,altFile$\;
}
\end{algorithm}

The algorithm first tries to locate the method by matching its signature (line 11), including the method name and parameters, but excluding the return type. If the signature match fails, it attempts to locate the method by the textual similarity of its body (line 14), considering a match valid if the similarity is at least 70\%. If the method still cannot be found, the search expands to other modified files (line 16), raising the similarity threshold to 75\% in this case. If the method is found in a different file—due to a method move or a file rename/move—the new file path is recorded in \texttt{altFile}. When the method is located through any of the above mechanisms (line 19), the algorithm checks whether the methods (from the current and parent commits) are identical, including their signatures, JavaDocs, annotations, or whether the file has been renamed or moved (in which case \texttt{altFile} should not be \texttt{NIL}) (line 20). If any change is detected—either in the method or in the file name or path—this change is recorded by invoking the \texttt{Add} procedure (line 21).

If the method is found in another file, a new DAG must be reconstructed for that file. To do so, the \texttt{HistoryFinder} procedure is recursively called (line 24), returning a sub-history that is appended to \( H \) (line 25). If the method is instead found within the same file (line 27), all immediate ancestor commit nodes are processed. In this case, all adjacent ancestor commit nodes of \( u \) that have not yet been visited are added to the visited set (line 30) and enqueued (line 33). The method name (line 31) and line number (line 32) are also updated to reflect the newly found version, so that it can be retrieved accurately when this node is dequeued later. Once all nodes have been processed, the resulting sequence of commit nodes captures the method’s change history in \(H\). Additionally, the furthest node dequeued from the queue $Q$ is added as the introduction commit, since this is the earliest known instance of the method.

Algorithm~\ref{alg:history-finder-similarity-matching} defines the utility procedures used to compare and match methods based on textual similarity. The procedure \texttt{SignatureMatch} attempts to locate the target method within a Java file at a specific commit. Line 2 parses all methods in the file using \texttt{JavaParser}, and lines 3 to 7 iterate over these methods, comparing each one’s signature with the target method. If a matching signature is found, the corresponding method is returned (line 5); otherwise, \texttt{NIL} is returned (line 8). The next procedure, \texttt{BodyMatch} (line 10), identifies a method based on the highest body similarity with the target method, provided it exceeds a specified similarity threshold. A threshold of 70\% is used when comparing methods within the same file, and 75\% when comparing across different files. Line 11 parses all Java methods in the file provided as input. The variable \texttt{bm} (line 12) tracks the best matching method, while \texttt{bs} (line 13) stores the corresponding similarity score. Line 14 iterates over all parsed methods, and the \texttt{Similarity} procedure (line 15) computes the similarity between the current method and the target method using the Jaro-Winkler algorithm~\cite{winkler90}. If the computed similarity meets the threshold and is greater than the current best score \texttt{bs}, both the score and the best matching method are updated (lines 16–19). Finally, the best matching method along with its similarity score is returned (line 21).

The procedure \texttt{AltFileMatch} attempts to locate the target method in all files that were modified in a given commit. It begins by listing all modified files using the \texttt{FindChangedFiles} procedure (line 24). Lines 25 to 27 initialize variables to track the best matching method, its similarity score, and the corresponding file path. Each modified file is then iterated (line 28), and the best matching method within each file is identified (line 29). If a method with a higher similarity score is found, the method information is updated accordingly (lines 30–34). Finally, the method with the highest similarity score and the new file path are returned (line 36). Similar to the \textit{CodeShovel} algorithm~\cite{grund_codeshovel_2021}, the first 100 methods from the corrected \textit{CodeShovel} oracle were used to train the algorithm and determine the optimal thresholds for body similarity within the original and other files.
